\documentclass[a4paper]{article} %

\date{} %

% Please, do not change any of the above lines

\usepackage{amssymb} % add other LaTeX packages you need here.



\begin{document}

\setcounter{page}{1} % Please, do not delete this line

% Your title
\title{Star-Free Modules}

% Authors
\author{Milan~Lekár, Jan~Paseka,  Richard Smolka\\ %
       Faculty of Science, Masaryk University\\
%Brno, Czech Republic
    \texttt{paseka@math.muni.cz} % Only one corresponding e-mail, please
       }%

\maketitle

% The abstract



Dynamic algebra provides an algebraic framework for expressing weakest preconditions of programs 
and has been widely used in formal program verification. When one wishes to verify \emph{quantum} 
programs, it is natural to extend dynamic algebra to settings where the underlying propositional 
logic is not classical but quantum---that is, where the lattice of propositions is an 
\emph{orthomodular lattice} rather than a Boolean algebra.

Orthomodular lattices arise naturally in quantum mechanics: the closed subspaces of any Hilbert 
space form a complete orthomodular lattice, the so-called Hilbert lattice. Crucially, 
orthomodular lattices are in general non-distributive, which distinguishes quantum logic from 
classical propositional logic and forces a careful re-examination of algebraic constructions 
that are routine in the Boolean setting.

A \emph{star-free dynamic orthomodular lattice} (introduced in \cite{Takagi}) is an orthomodular 
lattice equipped with a necessity operator $\Box(a, p)$---read as ``$p$ holds after executing 
program $a$''---for programs built from sequential composition, non-deterministic choice, 
and tests, but \emph{without} the Kleene star. The Kleene star is excluded because its treatment 
in the orthomodular/quantum setting raises substantial additional difficulties.

The main result of \cite{Takagi} is a \emph{Stone-type representation theorem}: every star-free 
dynamic orthomodular lattice embeds into its canonical extension, a complete star-free dynamic 
orthomodular lattice constructed from the proper filters of the original one. This mirrors the 
classical Stone representation of Boolean algebras by set algebras, and it establishes one half 
of a discrete duality between star-free dynamic orthomodular lattices and the corresponding 
Kripke-style frames (star-free dynamic orthomodular frames).

The present paper develops a \emph{module-theoretic} perspective on this structure. Rather than 
the necessity operator $\Box(a,p)$ of \cite{Takagi} (which corresponds to the \emph{weakest 
precondition}), we work with a dual action $\ell(a,p)$ that propagates information 
\emph{forward} along programs. The program terms form a term algebra $QA[\Pi, L]$ which, after 
quotienting by the natural behavioural equivalence, yields a unital involutive $m$-semilattice 
$S_{\mathbf{Dl}}$ acting on the orthomodular lattice $L$ (see \cite{BLP}). This algebraic structure is the 
module-theoretic analogue of the complex algebras studied in \cite{Takagi}.

A central subtlety---illustrated by a detailed example ---is that 
the involution on programs does not descend to the naive quotient by behavioural equivalence of 
forward actions alone. One must simultaneously track the behaviour of the dual program $a^*$, 
leading to the strengthened equivalence $\approx_{\mathbf{Dl}}$. This phenomenon has a direct 
counterpart in the frame semantics of \cite{Takagi}, where the self-adjointness axiom for 
test relations plays an analogous role.




\begin{thebibliography}{99}

\bibitem{BLP}
M.~Botur, J.~Paseka, M. Lekár,  
\newblock \emph{Foulis m-semilattices and their modules,} 
\newblock  In: {Proceedings of the 55th International Symposium on Multiple-Valued Logic (ISMVL 2025)}, pp. 196--201, IEEE (2025).
\bibitem{Takagi}
T.~Takagi.
\newblock \emph{A Stone-type representation theorem for star-free dynamic   orthomodular lattice,}
\newblock In: {Proceedings of the 56th International Symposium on Multiple-Valued Logic (ISMVL 2026)}, pp. 129--134, IEEE (2026). 

\end{thebibliography}

\end{document}
